% TEF16: a fixed-field GoldenFloat with a balanced-ternary exponent.
% v1: 2026-08-08. Accuracy re-measured independently; hardware measured on
% XC7A200T through the open flow. House style follows research/arxiv_submission.
\documentclass[11pt]{article}

\usepackage[utf8]{inputenc}
\usepackage[T1]{fontenc}
\usepackage{hyperref}
\usepackage{graphicx}
\usepackage{booktabs}
\usepackage{amsmath}
\usepackage{amssymb}
\usepackage{url}
\usepackage{geometry}
\usepackage{tabularx}
\usepackage{array}
\usepackage{enumitem}
\usepackage{amsthm}
\newtheorem{theorem}{Theorem}
\newtheorem{proposition}[theorem]{Proposition}
\newtheorem{corollary}{Corollary}
\theoremstyle{definition}
\newtheorem{definition}{Definition}

\geometry{a4paper, margin=1in}

\hypersetup{
  colorlinks=true,
  linkcolor=blue,
  urlcolor=blue,
  citecolor=blue
}

% The title names the category and the design point rather than an internal
% codename. tekum titles itself "balanced-ternary tapered-precision arithmetic";
% naming the opposite choice in the same breath is the positioning, and a reader
% who knows that family sees the contrast without needing the codename.
% The thesis is what the ladder is FOR: a reference whose error is predictable
% from its parameters and does not move with magnitude. Being fastest or most
% accurate at one point is not what a ruler needs to be.
\title{Ternary-Exponent Floats: a fixed-field ladder with no regime to decode \\
  \large A balanced-ternary exponent over a binary radix, from 4 to 1024 bits,
  measured against the published posit and takum codecs on open silicon}
\author{Dmitrii Vasilev\\
  \small ORCID 0009-0008-4294-6159 \\
  \small \texttt{admin@t27.ai}}
\date{9 August 2026}

\begin{document}
\maketitle

\begin{abstract}
A reference format is not the one that is fastest, or the most accurate at any one
magnitude. It is the one whose error you can predict from its parameters and which
does not change depending on where you look. We describe \textbf{Ternary Floats}
(TEF), a ladder built for that role: the exponent is a balanced-ternary integer of
$E_t$ trits, the mantissa is a constant $M$ bits, and the fields are fixed.

The foundation is a law rather than a table. Across every rung the mean relative
round-trip error is $\tfrac12\mathbb{E}[1/s]\,2^{-(M+1)}$
(Theorem~\ref{thm:law}), derived rather than fitted: the binade exponent cancels,
leaving a closed form in the mantissa width and the significand distribution
alone. Measured across eight rungs spanning 302 orders of magnitude the constant
holds to $\pm3\%$, and within each rung the error varies by at most $7\%$ from one
end of a 76-binade workload to the other.

The law's consequence is that a format's costs become commensurable. A mantissa
bit has a marginal area $\lambda$, so a variable-length exponent field costs a
measurable number of mantissa bits of silicon (Theorem~\ref{thm:convert}).
Synthesising the regime codecs of both tapered families under one harness on an
XC7A200T, net of harness: a unary regime costs $438$ LUTs and a length-prefixed
one $40$, against $0$ for a fixed field, which at the 16-bit class are $8.98$ and
$0.82$ mantissa bits. The first exceeds TEF16's entire mantissa; the second is
small. On that axis the argument is decisive against posit and weak against takum,
and we say so rather than averaging the two.

Synthesising the \emph{published} decoders rather than our models of them then
moves the main result somewhere else, and we had been measuring the smaller half.
\texttt{posit32\_decode} costs $480$ LUTs, as its structure predicts.
\texttt{takum32\_decode} evaluates $2^{f}$ from a $65536\times 48$-bit table and
costs $10{,}967$ LUTs and $84$ RAMB36 tiles --- $23\%$ of the device's block
memory --- because takum's value law is logarithmic and a linear datapath needs
the exponential. TEF's exponent path is wires. Theorem~\ref{thm:valuelaw} makes
this a second axis, orthogonal to tapering: what a format costs to \emph{read} is
set by whether its value law is $s\cdot 2^{e}$ or $2^{\ell}$, not by how it spends
precision. Unlike the accuracy result, Theorem~\ref{thm:floor} places no range
bound on it. One asymmetry is stated rather than left implicit: a logarithmic law
turns multiplication into addition, and a design that stays in the logarithmic
domain never pays this.

What remains against takum is accuracy, and Theorem~\ref{thm:floor} bounds how far
that claim may be taken. The exponent's encoding is a prefix code, so by Kraft no
format of unbounded range tapers more gently than one bit per doubling; takum's
regime is asymptotically optimal and nothing, TEF included, beats it beyond its
own range. Inside TEF16's range it is $3.70\times$ and $11.21\times$ more accurate
at mid and far magnitudes, and $3.70\times$ near unity where the unfilled rung
used to lose.

Turned around, the law is an instrument. The mantissa width a format
\emph{behaves} as having, measured band by band inside each format's own range,
resolves an 83-format catalogue into four staircase forms and only four: constant
(38 formats), wobble (3, at exactly $\log_2 r$ bits, reproducing IBM
hexadecimal's textbook figure), arithmetic ($5$; posit at exactly $2^{-es}$ per
binade, confirmed at three widths) and geometric ($8$; takum and tekum at one bit
per doubling).

In silicon on an entirely open flow: $12$, $50$, $212$, $1{,}477$ and $7{,}479$
LUTs at $161.1$, $153.2$, $131.7$, $83.3$ and $48.2$\,MHz for TEF4 through
TEF64, one cycle of latency, no DSP blocks.

We report what the measurements cost us. Three defects found while measuring
generalise and are described rather than quietly fixed. Four claims of our own
were falsified by the measurements meant to support them, including a predicted
area for TEF64 that came in $1.57\times$ low, and each is retracted in place.

Read together the results compose into a claim about the design space rather than
about one format. It has four axes; three admit answers that do not depend on the
workload. The radix of the scale is fixed at two unconditionally --- scaling is a
shift only for $r = 2^{k}$, and on that set the error constant is strictly
increasing, so the optimum is not the unconstrained extremum near $r = 1.9$ but
the best point that can be built. The radix of the exponent's encoding is fixed by
the fabric. Fixed fields or a taper is a genuine choice with exactly two options.
The exponent--mantissa split follows once a range is named. Name the fabric, the
range and the side of the dichotomy, and the format is determined. TEF is what
that returns for one triple, which makes it forced rather than designed --- and
carries no claim off those arguments.

\medskip
\noindent\textbf{Artifacts.} Oracle, RTL, testbenches and the exact commands are
public; every figure here regenerates from them.
\end{abstract}

\section{What this paper argues}
\label{sec:roadmap}

The results below were found in the order a measurement campaign finds things,
but they compose into a single claim, and it is easier to read them knowing it:

\begin{quote}
A floating-point format has four design axes. Three of them have answers that do
not depend on the workload. Name the fabric, the range the workload actually
visits, and whether it wants bounded range at constant precision or unbounded
range at decaying precision --- and the format is determined, not chosen.
\end{quote}

Section~\ref{sec:optimal} states that as Corollary~\ref{cor:determined} and shows
TEF is what it returns for one particular triple. Everything before it establishes
an axis:

\begin{itemize}
\item \textbf{The split between exponent and mantissa} is fixed once a range is
named, because the error depends on the mantissa width alone
(Section~\ref{sec:law}). The same law, inverted, becomes an instrument that reads
a format's behaviour back out of its round-trip error, and resolves a
83-format catalogue into four staircase forms.
\item \textbf{Fixed fields or a taper} is a genuine choice with exactly two
options, because a prefix code cannot give unbounded range and bounded precision
at once (Section~\ref{sec:codes}). That section also shows the taper's shape is
the growth rate of a codeword length, nothing more.
\item \textbf{The radix of the exponent's encoding} is decided by the fabric, not
by preference: packed into bits, a ternary exponent never spans more values per
bit than a binary one (Section~\ref{sec:whyternary}).
\item \textbf{The radix of the scale} is fixed at two unconditionally, and this is
the only axis with no free parameter at all. It is also the one where the optimum
is not the extremum (Section~\ref{sec:optimal}).
\end{itemize}

Two sections are about what things cost rather than what they are.
Section~\ref{sec:convert} makes area and precision commensurable, so a
variable-length exponent field can be priced in mantissa bits of silicon.
Section~\ref{sec:valuelaw} then finds that for a logarithmic format the field
layout is the smaller half of the cost, and that we had been measuring the wrong
half.

We also report what the campaign cost us. Four claims of our own were falsified by
the measurements meant to support them, and each is retracted where it was made
rather than quietly dropped.

\section{The format}\label{sec:format}

\begin{equation}
\text{TEF16}=[\,s\,|\,E{=}4\ \text{balanced-ternary trits}\,|\,M{=}11\ \text{bits}\,],
\qquad
v=(-1)^{s}\Bigl(1+\frac{M}{2^{9}}\Bigr)2^{e},
\quad e=\sum_{i} t_i 3^i \in [-40,40].
\end{equation}

Three properties follow from the layout rather than from tuning.

\paragraph{No regime decode.} The dominant cost in a tapered format is the
variable-length regime: it must be located, its length counted, and the
significand barrel-shifted into place. That cost is paid on any fabric, ternary
included. Fixed fields make it a bit slice.

\paragraph{The exponent is already ternary.} Four balanced-ternary trits give
$3^4=81$ exponent steps. On a ternary fabric the exponent add is native --- no
binary carry, no conversion between bases.

\paragraph{Precision does not taper.} Nine mantissa bits at every magnitude,
where a tapered format narrows towards the extremes. This is the mechanism behind
the accuracy result of Section~\ref{sec:accuracy}: the win appears where the
comparison format has spent its mantissa on range.


\section{Why a ladder, and why this one}\label{sec:ruler}

A format that is used as a reference has different requirements from one used as
a workhorse. A workhorse should be fast and accurate where the data is. A
reference has to be \emph{predictable}: given its parameters you should know its
error without measuring, and that error should not move when you look at a
different part of the range. A ruler whose divisions grow towards one end is not
a ruler.

Table~\ref{tab:law} tests both properties on this ladder.

\begin{table}[t]\centering
\caption{The precision law and the flatness, measured in exact rationals.
``Flatness'' is the ratio of the worst magnitude band to the best: $1.000$ would
mean the error does not depend on magnitude at all.}
\label{tab:law}
\begin{tabular}{lrrrrr}
\toprule
rung & $M$ & $2^{-(M+1)}$ & measured & ratio & flatness \\
\midrule
TEF8    & $4$    & $3.12\mathrm{e}{-2}$   & $1.22\mathrm{e}{-2}$   & $0.39$ & $1.000$ \\
TEF16   & $9$    & $9.77\mathrm{e}{-4}$   & $3.60\mathrm{e}{-4}$   & $0.37$ & $1.070$ \\
TEF32   & $25$   & $1.49\mathrm{e}{-8}$   & $5.49\mathrm{e}{-9}$   & $0.37$ & $1.070$ \\
TEF64   & $52$   & $1.11\mathrm{e}{-16}$  & $4.22\mathrm{e}{-17}$  & $0.38$ & $1.050$ \\
TEF128  & $115$  & $1.20\mathrm{e}{-35}$  & $4.44\mathrm{e}{-36}$  & $0.37$ & $1.070$ \\
TEF256  & $242$  & $7.07\mathrm{e}{-74}$  & $2.69\mathrm{e}{-74}$  & $0.38$ & $1.050$ \\
TEF512  & $497$  & $1.22\mathrm{e}{-150}$ & $4.51\mathrm{e}{-151}$ & $0.37$ & $1.070$ \\
TEF1024 & $1006$ & $7.29\mathrm{e}{-304}$ & $2.77\mathrm{e}{-304}$ & $0.38$ & $1.050$ \\
\bottomrule
\end{tabular}
\end{table}

\paragraph{The precision is a closed form, and it is derived rather than fitted.}
Write $u = 2^{-(M+1)}$ for the unit roundoff. Inside a binade the absolute
rounding error under round-to-nearest is uniform on $[-u\,2^{e}, +u\,2^{e}]$
while the value is $v = s\,2^{e}$ with significand $s \in [1,2)$, so the relative
error is $|U|/s$ with $|U|$ uniform on $[0,u]$ and \emph{independent of $e$}.
Hence
\begin{theorem}[Precision law]\label{thm:law}
For a format with a constant significand width of $M$ bits under
round-to-nearest, with significand $s$ and exponent $e$ independent,
\[
  \mathbb{E}\bigl[|\text{rel err}|\bigr] \;=\; \tfrac{1}{2}\,\mathbb{E}[1/s]\;\cdot\;2^{-(M+1)} ,
\]
independently of $e$.
\end{theorem}
The exponent drops out entirely --- that is the flatness of the next paragraph,
proved rather than observed. The constant is not universal: it is
$\tfrac{1}{2}\ln 2 = 0.3466$ for a significand uniform on $[1,2)$ and
$\tfrac{1}{2}\,(2\ln 2)^{-1} = 0.3607$ under a logarithmic (Benford) significand.
Our workload has $\mathbb{E}[1/s] = 0.7721$, predicting $0.3861$; the measured
figures in Table~\ref{tab:law} average $0.3756$ over eight rungs, with a spread
of $0.369$--$0.390$ that brackets it.

The practical consequence is the one a reference needs: given $M$ and the
significand distribution of a workload, the error of a rung nobody has built is
known before it is built.

\paragraph{The error does not move with magnitude, and cannot.} The derivation
above removes $e$ from the expression, so magnitude-independence is a property of
the layout rather than a happy measurement. Empirically flatness stays between
$1.05$ and $1.07$ --- at most a $7\%$ difference between the closest and furthest
parts of a workload spanning 76 binary decades. This is the property fixed fields
buy and the one a tapered format cannot have: spending mantissa bits on range is
exactly what makes precision depend on where you are.

\paragraph{What this does not claim.} Not that TEF is the most accurate format
at any given width --- Section~\ref{sec:field} shows posit16 and binary16 ahead
near unity, and Section~\ref{sec:field} shows the binary formats ahead outright at
32 bits. A reference does not need to win those comparisons. It needs to be the
thing you measure them against.


\section{The law is general, and it measures tapering}\label{sec:law}

Section~\ref{sec:ruler} derived the error of a fixed-field format. Nothing in that
derivation is specific to TEF, so it should hold for any format whose significand
width does not vary --- and should visibly fail for one whose width does. That
makes it an instrument rather than a description, and this section uses it as one.

\begin{definition}[Effective mantissa]
For a format and a magnitude band $B$, define
\[
  M_{\mathrm{eff}}(B) \;=\; -\log_2\!\left(\frac{2\,\mathbb{E}\bigl[|\text{rel err}|\;\big|\;B\bigr]}{\mathbb{E}[1/s]}\right) - 1 ,
\]
the mantissa width a fixed-field format would need to produce the error observed
in $B$.
\end{definition}

\begin{theorem}[Diagnostic]\label{thm:diag}
$M_{\mathrm{eff}}$ is constant across bands if and only if the format holds a
constant significand width over those bands and does not exhaust its range. For a
tapered format $M_{\mathrm{eff}}$ decreases with $|e|$, and the slope
$\mathrm{d}M_{\mathrm{eff}}/\mathrm{d}|e|$ is the taper rate in bits of
significand per binade.
\end{theorem}

Table~\ref{tab:meff} applies it. The measurement is a uniform significand on
$[1,2)$ so that $\mathbb{E}[1/s] = \ln 2$ exactly, removing the workload from the
question.

\begin{table}[t]\centering
\caption{Effective mantissa by magnitude band, and the fitted taper rate. Declared
$M$ is recovered to $0.01$ bits for every fixed-field format, which validates the
law; the tapered formats separate cleanly.}
\label{tab:meff}
\begin{tabular}{lrrrrrr}
\toprule
format & declared $M$ & $|e|\,0$--$6$ & $6$--$12$ & $12$--$20$ & $20$--$30$ & bits/binade \\
\midrule
TEF16   & $9$  & $8.99$  & $8.99$ & $9.00$ & $9.01$ & $+0.001$ \\
GF16     & $9$  & $8.99$  & $8.99$ & $9.00$ & $9.01$ & $+0.001$ \\
bfloat16 & $7$  & $7.04$  & $7.04$ & $7.00$ & $7.04$ & $-0.000$ \\
binary16 & $10$ & $10.01$ & $10.01$ & $7.41$ & $0.52$ & --- \\
posit16  & ---  & $10.72$ & $9.34$ & $7.48$ & $5.17$ & $\mathbf{-0.254}$ \\
takum16  & ---  & $9.62$  & $8.17$ & $7.34$ & $7.04$ & $\mathbf{-0.113}$ \\
\bottomrule
\end{tabular}
\end{table}

\paragraph{The law recovers what the formats declare.} TEF16 and GF16 return
$8.99$--$9.01$ against a declared 9; bfloat16 returns $7.00$--$7.04$ against a
declared 7. A law that reconstructs a format's own parameter to a hundredth of a
bit, from nothing but its round-trip error, is doing more than curve-fitting.

\paragraph{Tapering becomes a number.} posit16 sheds $0.254$ bits of significand
per binade ($R^2 = 0.999$) and takum16 appears to shed $0.113$ over the thirty
binades of this sweep --- but see Section~\ref{sec:taperlaw}: takum's taper is
logarithmic, so that figure is an artefact of the window rather than a rate. That is the
quantity a tapered format trades for range, and to our knowledge it is not usually
reported as a single figure. It also explains the accuracy tables directly ---
posit16 starts $1.7$ bits ahead of TEF16 and ends $3.8$ behind.

\paragraph{Running out of range is a different failure, and the instrument
separates it.} binary16 is a fixed-field format, yet its $M_{\mathrm{eff}}$
collapses from $10.01$ to $0.52$. It is not tapering: it is subnormalising and
then overflowing. A constant slope indicates a taper; a cliff indicates a range
limit. The two look alike in an accuracy table and do not look alike here.

\begin{corollary}
A format's position on the range--precision trade can be stated as a pair: the
constant $M_{\mathrm{eff}}$ it holds, and the number of binades over which it
holds it. For the TEF ladder those are $(M, 2\cdot 3^{E_t - 1})$ by construction;
for a tapered format neither is constant, which is why a ladder makes a better
ruler than a taper.
\end{corollary}


\subsection{The diagnostic across a catalogue}\label{sec:catalogue}

Applying Theorem~\ref{thm:diag} to every format for which a published oracle was
available --- 51 in all, spanning IEEE, bfloat, posit, takum, LNS, MXFP, VAX, IBM
hexadecimal, Cray, PDP-11, x87 and Microsoft binary, plus both GoldenFloat ladders
--- produced three observations we did not anticipate.

\paragraph{Declared widths are recovered across three orders of magnitude of
width.} The GF ladder returns $8.97$ at GF16 against a declared 9, $78.03$ at
GF128 against 78, and $631.99$ at GF1024 against 632. The same holds for
double-double ($106.68$), quad-double ($215.90$) and x87's 80-bit format
($63.04$). A relation that reconstructs a declared parameter from round-trip error
alone, at widths from 8 to 1006 bits, is doing something other than describing the
format it was derived on.

\paragraph{The taper rate is a constant of the family --- but only posit has a
rate.} posit16 sheds $0.254$ bits of significand per binade and posit32 sheds
$0.247$, against the $2^{-es} = 0.25$ that Theorem~\ref{thm:posit} below requires;
posit64, whose oracle carries the pre-standard $es = 3$, sheds $0.123$ against a
required $0.125$. Three formats, three exact confirmations. We previously reported
a comparable per-binade rate for takum, and that was a category error on our part:
takum's staircase is logarithmic, so fitting a line to it produces a number that
depends on the interval fitted and describes nothing. Its correct constant is one
bit per \emph{doubling} of $|e|$, and against that the four takum and three tekum
rungs measure $0.83$ to $1.09$.

\paragraph{Range exhaustion is distinguishable from tapering, and both from
neither.} The distinction is not optional: our own first pass placed the
measurement bands at fixed positions rather than inside each format's measured
range, and read the saturation of every narrow format as a taper --- classifying
TEF16, a fixed-field format, as geometrically tapered. Measuring each format's
usable range first and then placing the bands inside it removes the artefact
entirely. GF14 reads flat at $8.05$ across its range and collapses only below its
subnormal boundary, which is not a defect; the instrument says which of the two it
is, where an accuracy table alone does not.


\subsection{Tapering has a functional form, and it follows from the encoding}
\label{sec:taperlaw}

Section~\ref{sec:catalogue} reported a taper rate per family. That was half right,
and the half that was wrong is instructive: a rate presupposes a linear taper, and
only one of the two families has one.

\begin{theorem}[Posit taper, exact]\label{thm:posit}
For a posit of $n$ bits with exponent field $es$, the regime is a run of $k$ like
bits carrying a scale of $\mathrm{useed}^k = 2^{2^{es}k}$ and occupying $k+1$ bits.
Reading $k$ out of the encoder for $2^{e}$ confirms $k = \lfloor |e|/2^{es}\rfloor + 1$
for every $e$ tested. The significand therefore loses
\[
  \Delta M_{\text{posit}}(e) \;=\; \Bigl\lfloor \frac{|e|}{2^{es}} \Bigr\rfloor
  \quad\text{bits: a staircase, arithmetic in } |e| .
\]
\end{theorem}

The Posit Standard (2022) fixes $es = 2$ at every precision~\cite{gustafson}, so
the step is one bit per four binades at all widths --- which is why the loss
appeared as a family constant of $0.25$ bits per binade. Measured slopes of
$-0.2497$ and $-0.2505$ at posit16 and posit32 are the staircase seen through a
linear fit.

\begin{theorem}[Takum taper, exact]\label{thm:takum}
Takum spends a fixed 3-bit regime field holding $r$, after which the
characteristic occupies $r$ bits. Reading $r$ out of the encoder for $2^{e}$ gives
$r = \lfloor \log_2(|e|+1) \rfloor$ exactly, at 16, 32 and 64 bits alike. The
significand therefore loses
\[
  \Delta M_{\text{takum}}(e) \;=\; \bigl\lfloor \log_2(|e|+1) \bigr\rfloor
  \quad\text{bits: a staircase, geometric in } |e| .
\]
\end{theorem}

\begin{corollary}[The two families differ in kind, not degree]
Within its range a fixed-field format loses nothing; a posit loses one bit of
significand per $2^{es}$ binades without limit; a takum loses one bit per
\emph{doubling} of $|e|$. Tabulating $\Delta M$:

\begin{center}
\begin{tabular}{rrrr}
\toprule
$|e|$ & fixed-field & posit ($es{=}2$) & takum \\
\midrule
$30$      & $0$ & $7$    & $4$ \\
$100$     & $0$ & $25$   & $6$ \\
$1{,}000$   & $0$ & $250$  & $9$ \\
$10{,}000$  & $0$ & $2{,}500$ & $13$ \\
\bottomrule
\end{tabular}
\end{center}

At $|e| = 1{,}000$ a posit would need 250 bits of significand merely to pay for
its regime, which is why 16- and 32-bit posits cannot reach there at all. This is
the mechanism behind takum's stated improvement on posits~\cite{takum}, stated as
a quantity.
\end{corollary}

\begin{corollary}[A rate is not always the right summary]
Both laws are staircases, so a fitted slope is a summary and not the thing itself,
and it is a faithful summary only when the staircase is arithmetic. An earlier
draft of this paper reported $-0.117$ bits per binade for takum16 over
$|e| \in [0,30)$; the same format over $[0,60)$ gives a smaller figure for no
reason but the window, and the smooth-log fit of $c = 0.75$ understates the true
step of one bit per doubling for the same reason. Identify the form before
choosing a statistic --- and prefer reading the encoder to fitting its output.
\end{corollary}

\begin{corollary}[A three-way taxonomy, recovered without knowing the encoding]
The diagnostic separates formats by the \emph{form} of $M(e)$: constant
(fixed-field, TEF and GF by construction), linear in $|e|$ (posit), or
logarithmic in $|e|$ (takum). Which form fits is decided by the data, not by
reading the specification --- and the fitted coefficients then recover the
encoding parameters, $2^{-es}$ for posit and the characteristic width for takum.
\end{corollary}

What the literature describes qualitatively --- more precision near unity, less
at the extremes~\cite{gustafson} --- turns out to have an exact form recoverable
two ways that agree: from the encoder, by reading the regime field, and from
round-trip error alone, by the diagnostic of Theorem~\ref{thm:diag}. The second
route needs no access to the specification, which is what makes it useful on a
catalogue.


\subsection{Four staircase forms, and only four}\label{sec:taxonomy}

Classifying by the \emph{form} of the staircase rather than by a fitted slope
resolves the catalogue into four classes. Of the 83 entries, 11 are integer or
fixed-point formats with no exponent field; of the remaining 72, 18 span fewer
than six binades and cannot be classified from round-trip error at all, and one
requires probes wider than we generated.

\begin{table}[t]\centering
\caption{The catalogue by staircase form. The parameter is the taper constant in
the units natural to each form. Measured against each format's own usable range.}
\label{tab:taxonomy}
\begin{tabular}{llrl}
\toprule
form & $M_{\mathrm{eff}}$ against $|e|$ & count & members and constants \\
\midrule
constant   & flat                          & $38$ & IEEE, bfloat, VAX, Cray, x87, GF, TEF \\
wobble     & periodic, no trend            & $3$  & IBM hexadecimal, $\log_2 16 = 4$ bits \\
arithmetic & $-2^{-es}$ per binade         & $5$  & posit16/32 $0.25$, posit64 $0.125$ \\
geometric  & $-1$ per doubling             & $8$  & takum $\times 4$, tekum $\times 3$ \\
\bottomrule
\end{tabular}
\end{table}

The fourth class was not in our scheme until the catalogue produced it, and it has
a closed form.

\begin{theorem}[Precision wobble]
\label{thm:wobble}
Let a format scale by $r^{e}$ with a significand quantised uniformly over a
scale interval spanning a factor $r$. Its relative error varies by a factor of
$r$ across that interval, so $M_{\mathrm{eff}}$ oscillates with amplitude exactly
$\log_2 r$ bits and has no trend in $|e|$.
\end{theorem}

\begin{proof}
The quantisation step is constant across the interval and the relative error is
that step divided by the significand, which ranges over a factor $r$. Taking
$-\log_2$ of the ratio of the extremes gives $\log_2 r$. Since the interval
repeats identically at every $e$, the oscillation is periodic and carries no
trend. \qed
\end{proof}

Theorem~\ref{thm:wobble} reproduces a result sixty years older than this work.
IBM System/360 hexadecimal is documented as delivering between $4p-3$ and $4p$
significant bits, a spread of three~\cite{goldberg1991,ibmhfp}; the integer count
is the floor of the continuous quantity, so the theorem's $\log_2 16 = 4$ and the
textbook's $3$ are the same statement. Measured over eight sub-intervals ---
which by averaging within each bin predicts a reduced spread of $2.96$ rather than
the full $4$ --- \texttt{ibm\_hfp32} and \texttt{ibm\_hfp64} give $3.13$ and
$3.04$, while \texttt{binary32} and \texttt{gf32} give $0.85$ and $0.93$ against a
predicted $0.87$. Binary formats wobble by one bit; this is the familiar
$1$-bit wobble of every radix-2 float, and it is the smallest wobble any integer
radix admits. Theorem~\ref{thm:wobble} therefore supplies a second and independent
argument for TEF's binary scale, alongside the $0.331$ positions of
Theorem~\ref{thm:scaleradix} below: a radix-3 scale would wobble by $1.585$ bits
instead of $1$.

The taxonomy also settles why a format must choose between the first class and the
last two, and hence why this work is a ladder rather than a format.

\begin{theorem}[Range--precision dichotomy]
\label{thm:dichotomy}
Fix a width $N$. If a format's exponent takes unboundedly many values, then its
exponent codewords have unbounded length, and
$\lim_{|e| \to \infty} M_{\mathrm{eff}}(e) = 0$. Consequently no fixed-width
format has both unbounded range and precision bounded below by a positive
constant.
\end{theorem}

\begin{proof}
The exponent codewords form a prefix code over a binary alphabet, so by Kraft's
inequality $\sum_{e} 2^{-\ell(e)} \leq 1$. If the sum has infinitely many terms
then $\ell(e) \to \infty$. The payload available to the significand is
$N - 1 - \ell(e)$, so it falls to zero, and by Theorem~\ref{thm:law} so does
$M_{\mathrm{eff}}$. \qed
\end{proof}

\begin{corollary}[Why a ladder]
\label{cor:ladder}
Theorem~\ref{thm:dichotomy} makes constant precision and unbounded range
incompatible at any fixed width, so a format that wants both must obtain the
second by varying $N$. A ladder of fixed-field rungs does exactly this: each rung
holds its precision flat across its whole range, and the ladder as a whole covers
any range required. The tapered families make the opposite choice --- one rung,
unbounded range, precision decaying at the rate of their regime code, arithmetic
for posit and geometric for takum. Neither choice is better in the abstract; what
the dichotomy establishes is that they are the only two available.
\end{corollary}

Read against Corollary~\ref{cor:ladder}, the three tapered constants in
Table~\ref{tab:taxonomy} become a statement about codes rather than about
formats. A unary regime costs $\ell(e) \propto |e|$ and yields the arithmetic
staircase; a length-prefixed binary regime costs $\ell(e) \propto \log_2 |e|$ and
yields the geometric one; a fixed field costs $\ell(e) = \mathrm{const}$ and
yields no staircase at the price of bounded support. The staircase form is the
growth rate of the exponent's codeword length, and nothing else. This also names
the gap in the table: between $\ell \propto |e|$ and $\ell \propto \log_2|e|$ no
catalogued format sits. Section~\ref{sec:codes} builds the codes and asks what is
in the gap; the answer corrects a suggestion we made before building them.

\section{Accuracy}\label{sec:accuracy}

Mean relative error over an encode--decode round trip, 6{,}000 values spanning
$2^{-38}\ldots2^{38}$ with random sign, binned by binary-exponent magnitude
$|e|$. Table~\ref{tab:acc} and Figure~\ref{fig:acc}.

\begin{table}[t]\centering
\caption{Round-trip mean relative error. Bins are powers of two, not decades.}
\label{tab:acc}
\begin{tabular}{lrrrr}
\toprule
magnitude & GF16 ($\varphi$) & \textbf{TEF16} & tekum16 & TEF16 vs tekum16 \\
\midrule
$|e|<8$        & $3.43\mathrm{e}{-4}$              & $3.56\mathrm{e}{-4}$ & $3.27\mathrm{e}{-4}$ & $0.92\times$ (tie) \\
$|e|\,8$--$20$   & $3.57\mathrm{e}{-4}$              & $3.52\mathrm{e}{-4}$ & $1.00\mathrm{e}{-3}$ & $\mathbf{2.84\times}$ \\
$|e|\,20$--$38$  & $6.98\mathrm{e}{-3}$ (479 clip)   & $3.53\mathrm{e}{-4}$ & $1.95\mathrm{e}{-3}$ & $\mathbf{5.53\times}$ \\
\bottomrule
\end{tabular}
\end{table}

\begin{figure}[t]\centering
\includegraphics[width=0.82\textwidth]{tef_accuracy.pdf}
\caption{Accuracy by binary-exponent magnitude. GF16's 6-bit exponent clips 479
of 2{,}857 values in the far bin; TEF16's balanced-ternary exponent does not.}
\label{fig:acc}
\end{figure}

\paragraph{On the axis.} The bins are powers of two. Read as \emph{decades} the
far column is not a win at all: TEF16's exponent reaches $\pm40$ in powers of
two, roughly $\pm12$ decades, so beyond that it overflows while an unbounded
regime keeps working --- measured, 1{,}458 of the far-bin values clip under that
reading. We state the axis explicitly because an earlier internal note did not,
and the omission points a reader at the one interpretation under which the
result appears fabricated.

\section{The 16-bit field}\label{sec:field}

Comparing against one incumbent invites the objection that the incumbent was
chosen. Table~\ref{tab:field} runs every 16-bit-class format for which this
repository ships a published oracle~\cite{catalogue} through the same workload
and the same metric. Values that overflowed are counted separately rather than
folded into the mean, since a format should not look accurate for having
discarded the hard cases.

\begin{table}[t]\centering
\caption{Mean relative round-trip error by binary-exponent magnitude, with
overflow counts in parentheses. 6{,}000 values, $2^{-38}\ldots2^{38}$, random
sign, one seed.}
\label{tab:field}
\begin{tabular}{lrrr}
\toprule
format & $|e|<8$ & $|e|\,8$--$20$ & $|e|\,20$--$38$ \\
\midrule
\textbf{TEF16}    & $3.56\mathrm{e}{-4}$ & $3.52\mathrm{e}{-4}$ & $\mathbf{3.53\mathrm{e}{-4}}$ \\
GF16 ($\varphi$)   & $3.56\mathrm{e}{-4}$ & $3.52\mathrm{e}{-4}$ & $6.76\mathrm{e}{-3}$ (478) \\
takum16 / tekum16  & $3.27\mathrm{e}{-4}$ & $1.00\mathrm{e}{-3}$ & $1.95\mathrm{e}{-3}$ \\
posit16            & $\mathbf{1.36\mathrm{e}{-4}}$ & $8.31\mathrm{e}{-4}$ & $1.43\mathrm{e}{-2}$ \\
IEEE binary16      & $1.76\mathrm{e}{-4}$ & $1.28\mathrm{e}{-3}$ (299) & $1.57\mathrm{e}{-1}$ (2479) \\
bfloat16           & $1.45\mathrm{e}{-3}$ & $1.38\mathrm{e}{-3}$ & $1.41\mathrm{e}{-3}$ \\
LNS16              & $8.39\mathrm{e}{-4}$ (1324) & $7.63\mathrm{e}{-4}$ (1808) & $6.66\mathrm{e}{-4}$ (2849) \\
\bottomrule
\end{tabular}
\end{table}

\begin{table}[t]\centering
\caption{The whole ladder, measured in exact rationals. From TEF16 upwards the
error is flat in magnitude and its exponent roughly doubles per rung. TEF4 and
TEF8 are shown at their limits: the workload spans 76 binary decades and those
rungs hold 2 and 8, so most of it falls outside them --- reported rather than
omitted, since a ladder's lower rungs are where its range trade is visible.}
\label{tab:ladderacc}
\begin{tabular}{lrrrr}
\toprule
rung & decades & $|e|<8$ & $|e|\,8$--$20$ & $|e|\,20$--$38$ \\
\midrule
TEF4    & $2$      & \multicolumn{3}{c}{beyond range for this workload} \\
TEF8    & $8$      & $1.22\mathrm{e}{-2}$ & \multicolumn{2}{c}{beyond range} \\
TEF16   & $24$     & $3.45\mathrm{e}{-4}$   & $3.69\mathrm{e}{-4}$   & $3.66\mathrm{e}{-4}$ \\
TEF32   & $219$    & $5.27\mathrm{e}{-9}$   & $5.63\mathrm{e}{-9}$   & $5.58\mathrm{e}{-9}$ \\
TEF64   & $658$    & $4.33\mathrm{e}{-17}$  & $4.22\mathrm{e}{-17}$  & $4.12\mathrm{e}{-17}$ \\
TEF128  & $1{,}975$  & $4.25\mathrm{e}{-36}$  & $4.55\mathrm{e}{-36}$  & $4.51\mathrm{e}{-36}$ \\
TEF256  & $5{,}925$  & $2.76\mathrm{e}{-74}$  & $2.69\mathrm{e}{-74}$  & $2.63\mathrm{e}{-74}$ \\
TEF512  & $17{,}775$ & $4.32\mathrm{e}{-151}$ & $4.62\mathrm{e}{-151}$ & $4.58\mathrm{e}{-151}$ \\
TEF1024 & $53{,}326$ & $2.84\mathrm{e}{-304}$ & $2.77\mathrm{e}{-304}$ & $2.71\mathrm{e}{-304}$ \\
\bottomrule
\end{tabular}
\end{table}

\begin{figure}[t]\centering
\includegraphics[width=0.86\textwidth]{tef_ladder_acc.pdf}
\caption{The same data. Flatness across the three magnitude bands is the property
a fixed-field format is built for; the shaded rungs are those the workload
overruns.}
\label{fig:ladderacc}
\end{figure}


\begin{figure}[t]\centering
\includegraphics[width=0.92\textwidth]{tef_competition.pdf}
\caption{The same data. Scissors mark values that overflowed.}
\label{fig:field}
\end{figure}

Three readings, including the one that goes against us.

\paragraph{TEF16 is flat, and alone in being flat without clipping.} Its error
is $3.5\mathrm{e}{-4}$ at every magnitude in the workload and it discards
nothing. bfloat16 is also flat and clip-free, at roughly four times the error.
Every other entrant either degrades with magnitude or overflows: binary16 does
both, losing three orders of magnitude and discarding 2{,}479 values; GF16's
6-bit exponent clips 478; LNS16 clips throughout.

\paragraph{At 32 bits the binary formats win outright.} Running the same
workload on the 32-bit class: GF32 is flat at $3.4\mathrm{e}{-7}$ with no
clipping, but IEEE binary32 is flat at $2.1\mathrm{e}{-8}$, and posit32 and
takum32 are better still near unity ($2.1\mathrm{e}{-9}$ and
$4.9\mathrm{e}{-9}$). Uniformity stops being the deciding property once there are
enough bits that nothing clips anyway. The TEF argument is a 16-bit-class
argument, and we do not extend it.

\paragraph{Near unity we lose.} posit16 at $1.36\mathrm{e}{-4}$ and binary16 at
$1.76\mathrm{e}{-4}$ are more accurate than TEF16 there, by a factor of two and
a half and of two. A tapered format spends its bits where the values usually are,
and near unity that is the right bet. TEF16 buys uniformity instead, and the
trade only pays for a workload that leaves the neighbourhood of one.

\paragraph{takum16 and tekum16 are the same format here, exhaustively.} The two
oracles produce identical figures on every bin, so we checked the whole 16-bit
code space: all $65{,}536$ codes decode identically, with zero differences,
although the two implementations differ by 242 lines. The reconstruction our
tekum oracle implements \emph{is} takum16. Every comparison in this paper should
therefore be read as against takum16, and we relabel it as such. That is not a
weaker anchor --- takum is published, has a public VHDL codec, and is the parent
of the format we set out to compare against.

\section{Hardware}\label{sec:hardware}

Synthesis with Yosys~0.65~\cite{yosys}, place and route with
nextpnr-xilinx~(1743d0f)~\cite{nextpnr} on a
Xilinx XC7A200T, hard multipliers disabled so the figures describe fabric alone.
Bitstream generation uses Project X-Ray. No vendor tool is involved at any step,
which is what makes the numbers checkable by a reader without a licence.

\begin{table}[t]\centering
\caption{TEF multiplier, post-route, one harness across the whole lineup so the
rows are comparable. Latency one cycle, throughput one result per cycle.}
\label{tab:ladder}
\begin{tabular}{lrrr}
\toprule
rung & LUTs & DSP48 & $F_{\max}$ \\
\midrule
TEF4  & $12$      & $0$ & $161.11$\,MHz \\
TEF8  & $50$      & $0$ & $153.23$\,MHz \\
TEF16 & $212$     & $0$ & $131.73$\,MHz \\
TEF32  & $1{,}477$  & $0$ & $83.27$\,MHz \\
TEF64  & $6{,}140$  & $0$ & $53.26$\,MHz \\
TEF128 & $10{,}029$ & $0$ & $33.23$\,MHz \\
\bottomrule
\end{tabular}
\end{table}

\begin{table}[t]\centering
\caption{The specified ladder. $E_t$ is the trit budget, $M$ the mantissa width,
range in decades. Rungs above TEF128 exceed a single Artix-7 fabric and are
oracle- or ASIC-scale; they are specified and covered by the reference oracle,
not measured here.}
\label{tab:spec}
\begin{tabular}{lrrr@{\hskip 2em}lrrr}
\toprule
rung & $E_t$ & $M$ & decades & rung & $E_t$ & $M$ & decades \\
\midrule
TEF4  & $2$ & $1$  & $2.4$ & TEF64   & $7$  & $52$   & $658$ \\
TEF8  & $3$ & $4$  & $8$   & TEF128  & $8$  & $115$  & $1{,}975$ \\
TEF16 & $4$ & $9$  & $24$  & TEF256  & $9$  & $242$  & $5{,}925$ \\
TEF32 & $6$ & $25$ & $73$  & TEF512  & $10$ & $497$  & $17{,}775$ \\
       &     &      &       & TEF1024 & $11$ & $1006$ & $53{,}326$ \\
\bottomrule
\end{tabular}
\end{table}

\begin{figure}[t]\centering
\includegraphics[width=0.82\textwidth]{tef_ladder.pdf}
\caption{Area and achieved frequency across the ladder.}
\label{fig:ladder}
\end{figure}

For context rather than as a controlled comparison: a published Altera
\texttt{ALTFP\_MUL} on a Cyclone~IV reports 119--132\,MHz at 6--10 cycles of
latency, using 832--1041 logic elements \emph{and} 18 embedded multipliers, and
posit multiplier studies report 95--572 LUTs with 4.28--15.55\,ns of logic delay.
Different devices and different toolchains; we draw no equivalence from them
beyond the observation that TEF16's figures sit inside that envelope on all four
axes at once, on a fabric that gives the format no ternary advantage.

Isolating the pipeline register: the combinational multiplier closes at
$81.35$\,MHz, the two-stage version at $147.32$\,MHz. The cut is between the
significand product and the renormalisation --- a $10\times10$ multiply, then a
carry test, a saturating exponent add and a bit select.

\section{Two width defects, and why they generalise}\label{sec:width}

\subsection{Interface width dominated the arithmetic}

Our first synthesisable realisation declared every port and the significand
product 32 bits wide. Nothing in TEF16 is 32 bits: the mantissa field is 9, so
$1{+}M$ is 10 and their product is exactly 20; the exponent offset never exceeds
$\mathrm{OFFSET\_MAX}=80$, which is 7. Synthesis built a $32\times32$ multiplier
and a 32-bit comparison tree and charged for it (Figure~\ref{fig:width}):
1{,}179 LUTs, or three DSP48 blocks, against 219 LUTs and none once the buses
match the values they carry. The arithmetic is unchanged.

\begin{figure}[t]\centering
\includegraphics[width=0.82\textwidth]{tef_width.pdf}
\caption{The same arithmetic at three interface widths.}
\label{fig:width}
\end{figure}

\subsection{The same width decided correctness}

At TEF32 the significand product needs 52 bits. The 32-bit wire truncates, and
the module returns wrong results while appearing to support the rung, since its
own documentation names TEF32 as a supported configuration. Measured against a
64-bit reference: 1{,}995{,}730 mismatches in 2{,}128{,}964 combinations.

Deriving every width from the format parameters removes both problems at once. We
report this because the pattern is not specific to TEF: in a parametric
arithmetic core, a width that is merely \emph{large enough for the configuration
that was tested} is a silent-truncation trap at the next one, and it will not be
caught by a testbench written against the tested configuration.

\section{Equivalence}\label{sec:equiv}

Every change above is a claim that timing or width changed and arithmetic did
not. Each was checked rather than argued.

\begin{table}[t]\centering
\caption{Equivalence evidence. The reference for each rung is the module already
trusted for that rung.}
\label{tab:equiv}
\begin{tabular}{lrr}
\toprule
check & combinations / cycles & mismatches \\
\midrule
width-derived vs original, TEF16 (sweep)      & $321{,}156$   & $0$ \\
width-derived vs original, TEF4               & $374$         & $0$ \\
width-derived vs original, TEF8               & $3{,}716$     & $0$ \\
width-derived vs original, TEF16              & $77{,}444$    & $0$ \\
width-derived vs 64-bit reference, TEF32      & $300{,}000$   & $0$ \\
pipelined vs combinational                     & $199{,}994$   & $0$ \\
\midrule
original 32-bit module vs 64-bit ref, TEF32   & $2{,}128{,}964$ & $1{,}995{,}730$ \\
\bottomrule
\end{tabular}
\end{table}

The TEF16 sweep is exhaustive in the sense that matters here: the mantissa space
is covered in full at offset pairs exercising underflow, mid-range and
saturation, then the offsets are covered in full at mantissas that do and do not
carry --- every path through the carry, the saturation and the underflow clamp.
The last row is the point of the section: the width-derived module agrees with
the reference that is correct and disagrees with the one that truncates.


\section{What the law says about improving the ladder}\label{sec:design}

Theorem~\ref{thm:law} makes the design space arithmetic rather than a matter of
taste. Error scales as $2^{-(M+1)}$ and range as $3^{E_t}$, so at a fixed width
the only decision is how many bits the exponent field is allowed to consume ---
and that is decided by how efficiently $3^{E_t}$ values pack into
$\lceil E_t \log_2 3\rceil$ bits.

\begin{table}[t]\centering
\caption{Packing efficiency of the exponent field. The theoretical floor is
$\log_2 3 = 1.585$ bits per trit.}
\label{tab:pack}
\begin{tabular}{rrrrr}
\toprule
$E_t$ & values $3^{E_t}$ & bits & capacity & efficiency \\
\midrule
$3$  & $27$    & $5$  & $32$   & $84.4\%$ \\
$\mathbf{4}$  & $81$    & $7$  & $128$  & $\mathbf{63.3\%}$ \\
$5$  & $243$   & $8$  & $256$  & $\mathbf{94.9\%}$ \\
$6$  & $729$   & $10$ & $1024$ & $71.2\%$ \\
$\mathbf{7}$  & $2187$  & $12$ & $4096$ & $\mathbf{53.4\%}$ \\
$8$  & $6561$  & $13$ & $8192$ & $80.1\%$ \\
$10$ & $59049$ & $16$ & $65536$ & $90.1\%$ \\
\bottomrule
\end{tabular}
\end{table}

\paragraph{The width convention counts positions, and one rung breaks it.}
A rung's width is its position count with a trit occupying one position, so the
layout is $1 + E_t + M = N$. Three of the four specified rungs satisfy it exactly:
TEF4 is $1+2+1$, TEF8 is $1+3+4$, TEF32 is $1+6+25$. TEF16 is
$1+4+9 = 14$, two positions short of its name, and neither specification file
gives a reason. The likely history is visible in the parameters: TEF16 inherits
GF16's $\varphi$-optimal $M = 9$ and replaces GF16's six-bit exponent with four
trits, which buys range --- $3^4 = 81$ steps against $2^6 = 64$ --- and frees two
positions that were never re-spent. It is the one rung derived from its parent
rather than from the width rule.

This also settles what a rung costs on a binary fabric, which is a separate
question from its width: a trit stored as a two-bit code makes TEF16 eighteen
bits, and the offset stored as a plain binary integer makes it seventeen. Those
are fabric costs, not format widths, and confusing the two is what made the
ladder appear not to fit.

\paragraph{The slack is not confined to one rung.} Applying $1 + E_t + M = N$ to
every specified rung separates the two documents that define the ladder. The four
rungs of the machine-readable specification obey it at TEF4, TEF8 and TEF32;
the five upper rungs, which exist only in a research note, miss it by four to six
positions each, and the note and the specification disagree with each other about
TEF32.

\begin{table}[t]\centering
\caption{The width rule across the ladder. $k$ is the shortfall
$N - (1 + E_t + M)$. The reconciled column is $M = N - 1 - E_t$, which changes
nothing at the three rungs that already obey the rule.}
\label{tab:ladder-sweep}
\begin{tabular}{lrrrrrr}
\toprule
rung & source & $E_t$ & $M$ & $1{+}E_t{+}M$ & $k$ & reconciled $M$ \\
\midrule
TEF4    & spec & $2$  & $1$    & $4$    & $0$ & $1$ \\
TEF8    & spec & $3$  & $4$    & $8$    & $0$ & $4$ \\
TEF16   & spec & $4$  & $9$    & $14$   & $2$ & $\mathbf{11}$ \\
TEF32   & spec & $6$  & $25$   & $32$   & $0$ & $25$ \\
TEF32   & note & $5$  & $21$   & $27$   & $5$ & --- \\
TEF64   & note & $7$  & $52$   & $60$   & $4$ & $\mathbf{56}$ \\
TEF128  & note & $8$  & $115$  & $124$  & $4$ & $\mathbf{119}$ \\
TEF256  & note & $9$  & $242$  & $252$  & $4$ & $\mathbf{246}$ \\
TEF512  & note & $10$ & $497$  & $508$  & $4$ & $\mathbf{501}$ \\
TEF1024 & note & $11$ & $1006$ & $1018$ & $6$ & $\mathbf{1012}$ \\
\bottomrule
\end{tabular}
\end{table}

No single alternative convention explains the upper rungs either: reading the
exponent as $\lfloor E_t\log_2 3\rfloor$ bits fits TEF16, TEF64, TEF128 and
TEF1024 exactly but overshoots TEF256 and TEF512 by one and misses TEF32
entirely. We conclude that the upper ladder was tabulated rung by rung rather than
derived, and that reconciling it costs nothing at the rungs that are already
right.

\paragraph{What the two free positions were worth, and where they went.} Because
the exponent cancels in Theorem~\ref{thm:law}, the menu at sixteen positions could
be read off before measuring, and the measurement confirmed it to four significant
figures:

\begin{table}[t]\centering
\caption{The three ways to spend TEF16's two unallocated positions. Error and
range trade against each other exactly as Theorem~\ref{thm:law} predicts; nothing
in any row is worse than the unfilled rung.}
\label{tab:tef16-fill}
\begin{tabular}{lrrrr}
\toprule
variant & positions & mean rel.\ error & decades & vs.\ takum16 (far) \\
\midrule
unfilled, $E_t{=}4$, $M{=}9$  & $14$ & $3.37\mathrm{e}{-4}$ & $24$  & $5.46\times$ \\
$\mathbf{E_t{=}4}$, $\mathbf{M{=}11}$ (adopted) & $\mathbf{16}$ & $\mathbf{8.63\mathrm{e}{-5}}$ & $24$ & $\mathbf{22.68\times}$ \\
$E_t{=}5$, $M{=}10$           & $16$ & $1.69\mathrm{e}{-4}$ & $73$  & --- \\
$E_t{=}6$, $M{=}9$            & $16$ & $3.37\mathrm{e}{-4}$ & $\mathbf{219}$ & --- \\
\bottomrule
\end{tabular}
\end{table}

The last row is the law's own consistency check: adding two trits of exponent
while holding $M$ leaves the error bit-identical and multiplies the range
ninefold.

\paragraph{The positions were spent on mantissa, and the reason is a measurement.}
Corollary~\ref{cor:pair} forbids recommending a budget without naming a workload,
so we named one. At the 16-bit class the extra range is not consumed:
quantised-inference tensors span roughly $2^{-14}$ to $2^{14}$, and $E_t = 4$
already clips nothing at $\pm 40$ binades --- zero clips in 2000 values. $E_t = 6$
would have bought $\pm 364$ binades at bit-identical error, which is range nobody
spends at sixteen bits. So $M = 11$: the error falls by exactly $4.00$, and
against takum16 the rung moves from $0.92\times / 2.88\times / 5.49\times$ to
$3.70\times / 11.21\times / 22.68\times$, no longer losing the near bin. In
silicon it costs $443$ LUTs against $372$, a $19\%$ increase, at $136.44$\,MHz
against $131.73$ --- no frequency penalty.

All three ratios in this section and in Figure~\ref{fig:accuracy} come from one
script, \texttt{gen\_figures.py}, which computes them from the oracles at seed
$20260809$ rather than transcribing them. They were hardcoded until 2026-08-09 and
had drifted: the figure still carried $M = 9$ and the label \texttt{tekum16} after
the rung moved and the oracle was shown to decode identically to takum. A figure
that disagrees with the table it illustrates is the same defect class as a gate
that cannot go red.

Conformance was re-run in full rather than assumed: round-trip with sign
2000/2000, encoding monotone with zero inversions, $\pm$ symmetry with zero
mismatches in 500, both range boundaries decoding. A \texttt{width\_rule} test now
asserts $1 + E_t + M = N$ in the specification itself, so the position count is
checked rather than remembered. All nine rungs satisfy it.

\section{Why ternary, and exactly how much}\label{sec:whyternary}

The case for a ternary exponent is usually made by citing radix economy and left
there. We state it, then bound it in two directions --- what a binary fabric takes
back, and what the radix of the \emph{scale} is worth as distinct from the radix
of the \emph{exponent} --- because both bounds turn out to be decisive for the
design and neither is folklore.

\begin{theorem}[Radix economy]
\label{thm:economy}
Representing $R$ distinct values in radix $b$ costs $b\log_b R = b\ln R/\ln b$
state-slots. The factor $b/\ln b$ is minimised at $b = e$; among integers it is
$3/\ln 3 = 2.731$ against $2/\ln 2 = 2.885$, so a ternary field is $5.3\%$ more
economical than a binary one.
\end{theorem}

This is the classical argument~\cite{hayes2001,radixeconomy} and it is what
motivates a ternary exponent. It is also stated in a currency --- state-slots ---
that no binary fabric pays in. The next theorem says what happens when it does.

\begin{theorem}[No free range on a binary fabric]
\label{thm:nofree}
An exponent field of $E_t$ trits spans $3^{E_t}$ values. Stored as a binary
integer it occupies $\lceil E_t\log_2 3\rceil$ bits, a field spanning
$2^{\lceil E_t\log_2 3\rceil} \geq 3^{E_t}$ values. Hence on a binary fabric a
ternary exponent never spans more values per bit than a binary one, with equality
only at $E_t = 0$. Writing $\rho(E_t) = 3^{E_t}/2^{\lceil E_t\log_2 3\rceil}$ for
the utilisation, $\rho \in (1/2, 1]$ and $\rho$ does not converge.
\end{theorem}

\begin{table}[t]\centering
\caption{Utilisation $\rho(E_t)$ of the binary field holding an $E_t$-trit
exponent. The ladder's present budgets, $E_t \in \{4, 7\}$, are the two worst
entries in the range.}
\label{tab:rho}
\begin{tabular}{lrrrrrrrrrr}
\toprule
$E_t$ & 2 & 3 & 4 & 5 & 6 & 7 & 8 & 9 & 10 & 11 \\
\midrule
$\rho$ & $56.2$ & $84.4$ & $\mathbf{63.3}$ & $94.9$ & $71.2$ & $\mathbf{53.4}$ & $80.1$ & $60.1$ & $90.1$ & $67.6$ \\
\bottomrule
\end{tabular}
\end{table}

Run as an optimisation rather than an inequality, the same statement is sharper
than it looks. Asking which exponent encoding a binary fabric prefers at
$(N, R) = (16, 24)$, $(32, 219)$ and $(64, 658)$ returns the same mantissa width
for both --- $M = 10$, $23$ and $53$ respectively. The ternary encoding does not
lose on a binary fabric; at these widths the ceilings absorb the difference and it
exactly ties, which is the strongest form of ``never more values per bit''.

Theorem~\ref{thm:nofree} is the reason the ladder's rung widths must be read as
position counts and not bit counts, and it is where our own first reading of the
ladder went wrong: measuring what a trit costs when packed into bits and calling
the result the format's width confuses a fabric cost with a format property. The
independently published balanced-ternary float Ternary27~\cite{ternary27} counts
the same way --- 2 type trits, 1 sign trit, 5 exponent trits and 19 significand
trits, exactly 27 --- which is external corroboration that position counting is
the convention in this family rather than a reading of our own.

\begin{theorem}[Unallocated positions are a pure loss]
\label{thm:alloc}
Let a rung of $N$ positions carry a sign, $E_t$ exponent trits and $M$ mantissa
bits with $1 + E_t + M = N - k$, $k \geq 0$. Then relative to the rung with the
same $E_t$ and $M + k$, its expected relative error is larger by exactly $2^{k}$
and its dynamic range is identical.
\end{theorem}

\begin{proof}
By Theorem~\ref{thm:law} the expected relative error is
$\tfrac12 \mathbb{E}[1/s]\,2^{-(M+1)}$, in which the exponent does not appear;
replacing $M$ by $M+k$ multiplies it by $2^{-k}$. The dynamic range is
$3^{E_t}$ binades, a function of $E_t$ alone, so it is unchanged. Nothing else in
the layout depends on $M$. \qed
\end{proof}

Theorem~\ref{thm:alloc} converts the slack of Section~\ref{sec:design} into a
number without any measurement, and the measurement then confirms it. Against a
workload built in exact arithmetic --- necessary above TEF32, where a
double-precision probe measures its own resolution rather than the format's ---
the predicted and observed factors are:

\begin{table}[t]\centering
\caption{Unallocated positions per rung, the loss Theorem~\ref{thm:alloc}
predicts, and the loss measured against the reconciled rung
$M = N - 1 - E_t$. Exact-arithmetic workload, 800 probes per rung.}
\label{tab:alloc}
\begin{tabular}{lrrrr}
\toprule
rung & $k$ & predicted & measured \\
\midrule
TEF16  & $2$ & $4\times$  & $4.03\times$  \\
TEF64  & $4$ & $16\times$ & $15.43\times$ \\
TEF128 & $4$ & $16\times$ & $15.99\times$ \\
TEF256 & $4$ & $16\times$ & $15.24\times$ \\
\bottomrule
\end{tabular}
\end{table}

The final question is one the ladder answers implicitly and, as far as we can
find, no one has stated: TEF encodes its exponent in radix 3 but scales by
$2^{e}$, whereas Ternary27 scales by $3^{e}$. These are different choices and
they are separable.

\begin{theorem}[The scale radix, and why 2 beats 3]
\label{thm:scaleradix}
Let a format scale by $r^{e}$ with a normalised significand uniformly quantised
over $[1, r)$ by $M$ bits. Then its expected relative error on a log-uniform
workload is $\kappa(r)\,2^{-M}$ with
\[
\kappa(r) \;=\; \frac{(r-1)^{2}}{r\ln r},
\qquad \kappa(2) = 0.72135,\quad \kappa(3) = 1.21365 .
\]
Holding the field layout fixed, moving the scale from $r=2$ to $r=3$ therefore
multiplies the error by $\kappa(3)/\kappa(2) = 1.6825$ while multiplying the
dynamic range by $\log_2 3 = 1.5850$. Measured in positions, it costs
$\log_2 1.6825 = 0.7506$ of precision and buys $\log_3 1.5850 = 0.4192$ of range:
a net loss of $0.3314$ positions per number. A power-of-two scale strictly
dominates a power-of-three scale at equal width.
\end{theorem}

\begin{proof}
The quantisation step over $[1,r)$ is $(r-1)2^{-M}$ and the mean absolute
rounding error is half of it. On a log-uniform workload the significand has
density $1/(s\ln r)$ on $[1,r)$, so
$\mathbb{E}[1/s] = \frac{1}{\ln r}\int_1^{r} s^{-2}\,ds = \frac{1-1/r}{\ln r}$.
Multiplying gives $\tfrac12 (r-1)2^{-M}\cdot\frac{1-1/r}{\ln r}$, and
$\tfrac12(r-1)(1-1/r) = (r-1)^2/2r$. Absorbing the factor $\tfrac12$ into the
$2^{-(M+1)}$ convention of Theorem~\ref{thm:law} yields $\kappa$ as stated. For
the range, $E_t$ trits span $3^{E_t}$ steps of $\log_2 r$ binades each. One
mantissa bit halves the error and one exponent trit triples the range, which
fixes the exchange rate between the two currencies and gives the position
figures. \qed
\end{proof}

Direct measurement at two rungs, encoding and decoding through both scales in
exact arithmetic, gives error ratios of $1.675$ at $(E_t,M) = (4,11)$ and $1.690$
at $(6,25)$ against the predicted $1.6825$, and range ratios of $1.583$ and
$1.585$ against the predicted $1.5850$.

\begin{corollary}[What TEF is, precisely]
TEF is not a ternary float in the sense of Ternary27; it is a binary-radix float
whose exponent is \emph{encoded} in balanced ternary. Theorem~\ref{thm:scaleradix}
says that this is the better of the two available choices by $0.331$ positions per
number, and Theorem~\ref{thm:nofree} says that the ternary encoding pays only on a
ternary fabric. The two together delimit the format exactly: the radix choice is
right on any fabric, and the encoding choice is right on one.
\end{corollary}


\section{The regime field is a universal code}
\label{sec:codes}

Theorem~\ref{thm:dichotomy} identifies the exponent's encoding as a prefix code.
That identification is worth making explicit, because it places the taxonomy of
Section~\ref{sec:taxonomy} inside a hierarchy that predates every format in the
catalogue by decades~\cite{elias1975}.

\begin{theorem}[The staircase is the codeword-length function]
\label{thm:code}
Let a format of $N$ positions encode its exponent with a prefix code assigning
$e$ a codeword of length $\ell(e)$. Then
$M_{\mathrm{eff}}(e) = N - 1 - \ell(e)$, and the staircase form is the growth
rate of $\ell$ and nothing else.
\end{theorem}

The catalogue's three tapering classes are then three points in the classical
hierarchy: a unary regime, $\ell \propto |e|$, gives posit's arithmetic
staircase; a length-prefixed binary regime, $\ell \propto \log_2|e|$, gives
takum's geometric one; a fixed field, $\ell = \mathrm{const}$, gives no staircase
at the price of bounded support. Elias' $\gamma$, $\delta$ and $\omega$ codes
occupy the logarithmic class at successively finer asymptotics.

\begin{theorem}[Logarithmic is the floor]
\label{thm:floor}
For any prefix code on the integers, $\ell(e) \geq \log_2|e|$ for infinitely many
$e$. Hence no format of unbounded range tapers more gently than one bit per
doubling of $|e|$, asymptotically.
\end{theorem}

\begin{proof}
Kraft's inequality gives $\sum_e 2^{-\ell(e)} \leq 1$. If $\ell(e) < \log_2|e|$
for all but finitely many $e$, the tail of the sum exceeds
$\sum_e 1/|e|$, which diverges. \qed
\end{proof}

Theorem~\ref{thm:floor} is worth stating plainly because it bounds this work as
much as its competitors: takum's regime is asymptotically optimal, and no format
--- TEF included --- can beat it over unbounded range. What TEF does instead is
decline the unbounded range, which Theorem~\ref{thm:dichotomy} shows is the only
other option. The measured advantages of Section~\ref{sec:accuracy} are therefore
advantages \emph{inside TEF's range}, which is where the taper has already cost
takum the bits, and they should never be read as a claim beyond it.

\begin{theorem}[The regime's radix does not matter]
\label{thm:radixinv}
A regime spending one digit per multiplication of $|e|$ by $r$, over an alphabet
of $r$ symbols, costs $\log_2 r \cdot \log_r |e| = \log_2 |e|$ bits, independent
of $r$. The staircase form is invariant to the regime's radix.
\end{theorem}

Theorem~\ref{thm:radixinv} retracts a construction we proposed in an earlier draft
of this work. Having observed the gap between the arithmetic and geometric
classes, we suggested a regime spending one trit per tripling of $|e|$ as the
natural way for a ternary format to occupy it. It does not occupy it: building
both codes and normalising each to satisfy Kraft, the two codeword-length
functions differ by an additive constant of exactly $1$ at every $|e|$ from $1$ to
$4096$. One trit per tripling is takum's staircase with a different constant, not
a third form. A ternary fabric changes what the regime costs to decode; it does
not change what it costs to encode.

\paragraph{The gap is real, occupiable, and empty for a reason we can only
conjecture.} Growth rates strictly between $\log_2|e|$ and $|e|$ do exist and are
Kraft-feasible --- $\ell(e) = \lceil\sqrt{|e|}\rceil$ is one --- and such a code
is not dominated. At $N = 16$, normalised so that all three are legal prefix
codes:

\begin{table}[t]\centering
\caption{$M_{\mathrm{eff}}$ by $|e|$ and usable range for three regime codes at
$N = 16$, each normalised to satisfy Kraft. The square-root code is a genuine
Pareto point: more precision than the logarithmic code throughout the core, more
range than the unary one.}
\label{tab:codes}
\begin{tabular}{lrrrrrr}
\toprule
regime & $|e|{=}1$ & $4$ & $16$ & $64$ & $256$ & binades \\
\midrule
unary, $\ell \propto |e|$ (posit class)      & $11$ & $10$ & $7$ & $-5$ & $-53$ & $43$ \\
$\ell = \lceil\sqrt{|e|}\rceil$              & $11$ & $10$ & $8$ & $4$  & $-4$  & $121$ \\
$\ell \propto \log_2|e|$ (takum class)       & $9$  & $7$  & $5$ & $3$  & $1$   & $511$ \\
\bottomrule
\end{tabular}
\end{table}

For reference, a fixed field reaching the same 121 binades as the square-root code
spends 8 bits on the exponent and holds $7$ bits flat everywhere --- better than
the square-root code beyond $|e| \approx 30$ and worse below it, which is the
taper trade in miniature.

So the gap is not empty because its occupants are dominated; they are not. Our
first explanation was decode cost, and building the hardware refutes it.

\subsection{What a regime costs in silicon}\label{sec:regimecost}

We wrote three regime codecs to one interface --- a 16-bit word in, an exponent
and a mantissa shift out, and the mirror encoder --- and synthesised each under a
single harness on XC7A200T with Yosys and nextpnr-xilinx, hard multipliers
disabled. These are structural models of the three classes, not reproductions of
the published posit or takum codecs; what they compare is a run-scanned regime
against a length-prefixed one, which is the structural difference the taxonomy
names.

\begin{table}[t]\centering
\caption{Regime codec cost by class, post-route on XC7A200T under one harness,
no DSP. Figures are net of the harness, which costs $24$ LUTs and was synthesised
separately and subtracted; the round-trip column is decode plus encode.}
\label{tab:regimecost}
\begin{tabular}{lrrrrr}
\toprule
class & \multicolumn{2}{c}{decode} & \multicolumn{2}{c}{encode} & round trip \\
      & LUTs & $F_{\max}$ & LUTs & $F_{\max}$ & LUTs \\
\midrule
unary, $\ell \propto |e|$ (posit)          & $379$ & $36.8$\,MHz  & $59$  & $211.2$\,MHz & $438$ \\
$\ell \propto \log_2|e|$ (takum)           & $\mathbf{6}$ & $\mathbf{805.8}$\,MHz & $\mathbf{34}$ & $\mathbf{285.1}$\,MHz & $\mathbf{40}$ \\
$\ell = \lceil\sqrt{|e|}\rceil$            & $393$ & $33.1$\,MHz  & $123$ & $96.4$\,MHz  & $516$ \\
\midrule
fixed field (GF, TEF)                     & $0$   & ---          & $0$   & ---          & $\mathbf{0}$ \\
\bottomrule
\end{tabular}
\end{table}

The conjecture said the square-root code would be too expensive to be worth
building. It is not: it costs $18\%$ more than the unary regime round trip, and
the unary regime ships in production. What the measurement does show is that the
expense lies somewhere else than we assumed, and this has a closed form.

\begin{theorem}[Decode cost is set by the scan, not by the staircase]
\label{thm:scan}
A regime whose codeword length is delimited by a run must examine up to $N-1$
positions to decode, giving cost $\Theta(N)$. A regime whose length is carried in
a fixed field of $b$ bits requires no scan and a shift of at most $2^{b}-1$,
giving cost $\Theta(1)$ in $N$. The two costs are independent of the growth rate
of $\ell$, and therefore of the staircase form.
\end{theorem}

Theorem~\ref{thm:scan} predicts the table exactly. The unary and square-root codes
share the same 15-position scan and cost $403$ and $417$ LUTs, a difference of
$3.5\%$ that is the squarer and nothing else; the length-prefixed code has no scan
and costs $30$. The square-root code is geometrically intermediate and structurally
identical to the unary one, which is the point: \emph{the staircase form and the
decode cost are independent axes}, and a taxonomy of one does not predict the
other.

\begin{corollary}[The logarithmic regime is not a trade]
\label{cor:takumwins}
Against the unary regime the length-prefixed one is better on every axis measured:
asymptotically optimal by Theorem~\ref{thm:floor}, $11\times$ smaller round trip,
and $22\times$ faster to decode. The published bounded-regime posit
variants~\cite{positgap}, which cap the regime field precisely to avoid the scan,
are the same conclusion reached from the other direction.
\end{corollary}

\begin{corollary}[The measured price of unbounded range]
\label{cor:price}
A fixed field has no regime codec at all. The dichotomy of
Theorem~\ref{thm:dichotomy} therefore has a hardware statement: unbounded range
costs $40$ LUTs at best and $438$ at worst on this device, against $0$ for a
format that declines it. Section~\ref{sec:convert} converts those LUTs into the
same currency as the accuracy results, which is the only way to say whether they
matter.
\end{corollary}

We are left without an explanation for the empty gap. It is not domination, since
the square-root code is Pareto-undominated; and it is not cost, since it is within
$16\%$ of a regime that ships. Both of our explanations are now measured and both
are wrong, and we record that rather than replace it with a third.


\section{Converting silicon into precision}\label{sec:convert}

Corollary~\ref{cor:price} gives a regime's cost in LUTs and Section~\ref{sec:accuracy}
gives an accuracy ratio, and the two cannot be compared until they are in one
currency. They can be: at a fixed class, a mantissa bit has a marginal area, and
by Theorem~\ref{thm:law} it has a known value in accuracy.

\begin{theorem}[Area and precision are commensurable]
\label{thm:convert}
Let $\lambda = \mathrm{d}A/\mathrm{d}M$ be the marginal area of a mantissa bit at
a given class. A regime codec costing $C$ LUTs is then equivalent to $C/\lambda$
mantissa bits of silicon, and by Theorem~\ref{thm:law} to a factor
$2^{C/\lambda}$ in accuracy at equal area.
\end{theorem}

Measuring $\lambda$ directly --- the same width-derived multiplier synthesised at
$M = 7, 9, \ldots, 25$ under one harness --- gives $285$, $372$, $443$, $567$,
$696$, $870$, $1238$ and $1683$ LUTs. The marginal cost is not constant: it rises
from $35.5$ LUTs per bit between $M{=}9$ and $M{=}11$ to $111.2$ between $M{=}21$
and $M{=}25$. Taking the local value at each class, $\lambda = 48.8$ at the
16-bit class and $111.2$ at the 32-bit one:

\begin{table}[t]\centering
\caption{Regime area converted into mantissa bits through
Theorem~\ref{thm:convert}. The unary regime costs, in silicon alone, more than
TEF16's entire mantissa.}
\label{tab:convert}
\begin{tabular}{lrrr}
\toprule
regime & LUTs & bits at 16-bit class & bits at 32-bit class \\
\midrule
unary (posit class)          & $438$ & $\mathbf{8.98}$ & $3.94$ \\
length-prefixed (takum class) & $40$  & $0.82$ & $0.36$ \\
\bottomrule
\end{tabular}
\end{table}


\subsection{How area grows along the ladder}\label{sec:arealaw}

Theorem~\ref{thm:convert} needs $\lambda$, and $\lambda$ is a derivative of an
area law we can state. Fitting $A(M) = cM^{\alpha}$ to the eight rungs from
$M{=}7$ to $M{=}25$ gives $\alpha = 1.406$ at $R^2 = 0.984$, and combined with
Theorem~\ref{thm:law} it predicts a scaling relation:

\begin{proposition}[Error decays as a stretched exponential in area]
\label{prop:stretched}
With $\varepsilon \propto 2^{-M}$ and $A \propto M^{\alpha}$,
$\ln(1/\varepsilon) \propto A^{1/\alpha}$.
\end{proposition}

Measured over $M \in [7,25]$ --- a sixfold range of area and a $250{,}000$-fold
range of error --- the constant of proportionality holds to within $27\%$ and
drifts monotonically downward, so the relation is an approximation and we label
it a proposition rather than a theorem. The drift has a cause, and finding it
falsified a prediction we registered before synthesising.

\paragraph{The exponent is not constant, and the fit does not extrapolate.} We
predicted from the $[7,25]$ fit that the reconciled TEF64, at $M = 56$, would
cost $4{,}762$ LUTs. It costs $7{,}479$ at $48.20$\,MHz --- the prediction is low
by $1.57\times$. Reading $\alpha$ locally rather than globally shows why: it rises
monotonically, from $1.06$ between $M{=}7$ and $M{=}9$ to $1.78$ between $M{=}15$
and $M{=}17$ and $1.85$ between $M{=}25$ and $M{=}56$.

\begin{theorem}[The ladder's area law]
\label{thm:alpha}
A rung's area is a fixed overhead --- registers, the exponent adder, the
renormalisation --- plus a significand multiplier quadratic in $M$:
\[
A(M) \;=\; f + cM^{2}.
\]
Consequently the effective power-law exponent
$\alpha(M) = \mathrm{d}\ln A/\mathrm{d}\ln M$ is not a constant but rises
monotonically toward $2$, and the marginal cost of a mantissa bit,
$\lambda = \mathrm{d}A/\mathrm{d}M = 2cM$, is \emph{linear} in $M$ rather than
fixed.
\end{theorem}

Synthesising the same width-derived multiplier at fourteen mantissa widths from
$M{=}7$ to $M{=}90$ --- $285$, $372$, $443$, $567$, $696$, $870$, $1238$, $1683$,
$2601$, $3990$, $5934$, $7479$, $12232$ and $19980$ LUTs --- and fitting both
parameters gives $f = 141$ LUTs and $c = 2.4455$ LUTs per bit squared at
$R^2 = 0.99963$, with a maximum deviation of $9.9\%$ across a seventyfold range of
area. Read locally, $\alpha$ climbs from $1.06$ between $M{=}7$ and $M{=}9$ to
$2.20$ between $M{=}56$ and $M{=}70$, which is Theorem~\ref{thm:alpha} rather than
noise. On the registered prediction the structural model is right to $+4.4\%$
where the single power law was low by $36.3\%$; we report both, since the power
law was ours and we published the prediction before synthesising.

\begin{theorem}[A regime's worth in precision falls as $1/M$]
\label{thm:oneoverm}
By Theorems~\ref{thm:convert} and~\ref{thm:alpha}, a regime codec of fixed area
$C$ is worth
\[
\frac{C}{\lambda} \;=\; \frac{C}{2cM}
\]
mantissa bits, inversely proportional to the mantissa width. It is worth at least
one mantissa bit exactly while $M \leq C/2c$.
\end{theorem}

\begin{table}[t]\centering
\caption{The same regime codecs valued at five mantissa widths through
Theorem~\ref{thm:oneoverm}, with $c = 2.4455$ measured on XC7A200T. The unary
regime's silicon exceeds an entire TEF16 mantissa; the length-prefixed regime is
below one bit everywhere above the 8-bit class.}
\label{tab:oneoverm}
\begin{tabular}{rrrr}
\toprule
$M$ & $\lambda$ (LUT/bit) & unary, $438$ LUT & length-prefixed, $40$ LUT \\
\midrule
$9$  & $44.0$  & $\mathbf{9.95}$ & $0.91$ \\
$11$ & $53.8$  & $8.14$ & $0.74$ \\
$25$ & $122.3$ & $3.58$ & $0.33$ \\
$56$ & $273.9$ & $1.60$ & $0.15$ \\
$90$ & $440.2$ & $1.00$ & $0.09$ \\
\bottomrule
\end{tabular}
\end{table}

\begin{corollary}[Break-even widths]
\label{cor:breakeven}
With $c = 2.4455$, the unary regime is worth at least one mantissa bit of silicon
up to $M = 89.6$ --- that is, through the entire 128-bit class --- while the
length-prefixed regime reaches one bit only up to $M = 8.2$. takum's regime
therefore costs less than a single mantissa bit of silicon at every rung above
TEF8, and posit's costs more than one at every rung below TEF256.
\end{corollary}

Corollary~\ref{cor:breakeven} is the exact form of what
Corollary~\ref{cor:narrow} below could only say qualitatively, and it is the sharpest
statement we can make about where a fixed field pays for itself in area: against
posit, almost everywhere; against takum, only at the very bottom of the ladder.

\begin{corollary}[Range is paid for three times, and they are commensurable]
\label{cor:three}
A tapered format pays for unbounded range in three currencies: word positions,
in the regime field; silicon, in the codec; and precision at the extremes, in the
taper. Theorem~\ref{thm:convert} and Theorem~\ref{thm:law} put all three in
mantissa bits, so they can be added.
\end{corollary}

Applied honestly, Corollary~\ref{cor:three} separates the two tapered families
rather than lumping them. Against posit the silicon term alone is $8.98$ bits at
the 16-bit class --- more than TEF16's entire $9$-bit mantissa, so the area a
posit implementation spends decoding its regime would, spent on significand
instead, roughly double the precision budget of a fixed-field format at the same
class. Against takum the silicon term is $0.82$ bits, which is small. The
length-prefixed regime is cheap, and any argument for a fixed field made on area
grounds is an argument against posit and not against takum. Against takum what
remains is the accuracy result of Section~\ref{sec:accuracy}, bounded as
Theorem~\ref{thm:floor} requires to TEF's own range.

\begin{corollary}[The area argument is a narrow-width argument]
\label{cor:narrow}
Since $\lambda$ grows superlinearly in $M$ while a regime codec's cost is roughly
fixed, the same codec converts to fewer mantissa bits as the class widens:
$8.98$ bits at 16 becomes $3.94$ at 32 for the unary regime. The fixed-field area
advantage is therefore strongest where the words are narrowest, which is the
regime in which quantised inference operates and not the one in which numerical
linear algebra does.
\end{corollary}


\section{The value law costs more than the field layout}
\label{sec:valuelaw}

Sections~\ref{sec:regimecost} and~\ref{sec:convert} priced the regime as a field
to be extracted, and synthesising the published decoders shows that this is the
smaller half of the question. We take the takum and posit decoders from an
existing open verification suite --- Verilog, golden-checked against an mpmath
oracle --- and put them through the same harness on XC7A200T.

\begin{table}[t]\centering
\caption{Published 32-bit decoders, post-route on XC7A200T under one harness, no
DSP, converting to binary32. These are the real codecs rather than the structural
models of Table~\ref{tab:regimecost}.}
\label{tab:realcodecs}
\begin{tabular}{lrrrl}
\toprule
decoder & LUTs & RAMB36 & $F_{\max}$ & what dominates \\
\midrule
\texttt{posit32\_decode}          & $480$      & $0$  & $49.0$\,MHz & leading-zero count, barrel shift \\
\texttt{takum32\_decode} as published & $76{,}821$ & $0$  & ---         & the $2^{f}$ table, built from logic \\
\texttt{takum32\_decode} + clocked read & $10{,}967$ & $\mathbf{84}$ & $11.4$\,MHz & the same table, in block memory \\
TEF32, exponent path             & $0$        & $0$  & ---         & the field is read directly \\
\bottomrule
\end{tabular}
\end{table}

The takum row is given twice because one number alone would mislead. Its decoder
evaluates $2^{f_{\mathrm{hi}}/2^{16}}$ from a $65536 \times 48$-bit table, then
applies a Taylor correction through a $108$-bit and an $80$-bit multiply. As
published the table read is combinational, which no synthesiser can map to block
memory, so it was built from logic at $76{,}821$ LUTs. Registering that read and
marking it \texttt{ram\_style="block"} --- one cycle of added latency, and a
pipelined variant of the published module rather than the module itself --- puts
the table where it belongs: $10{,}967$ LUTs and $84$ RAMB36 tiles, which is
$3.1$\,Mbit and $23\%$ of this device's entire block memory. We predicted $85$
tiles from the table geometry before synthesising and measured $84$. Neither
number is a property of takum the format; both are properties of decoding a
logarithmic value law into a linear datapath, and neither is the regime field.

\begin{theorem}[Linear and logarithmic value laws]
\label{thm:valuelaw}
Let a format's value law be $v = s \cdot 2^{e}$ with $s$ and $e$ read from fields
(\emph{linear}), or $v = \pm 2^{\ell}$ for a decoded real $\ell$
(\emph{logarithmic}). A linear law decodes with a shift, cost independent of the
mantissa width. A logarithmic law requires evaluating an exponential, whose cost
is a table or an iterative kernel and is independent of the field layout ---
therefore independent of the staircase form of Section~\ref{sec:taxonomy}.
\end{theorem}

Theorem~\ref{thm:valuelaw} is a second axis, orthogonal to the first. The
staircase taxonomy classifies how a format spends \emph{precision}; the value law
classifies what it costs to \emph{read}. posit is tapered and linear; takum is
tapered and logarithmic; LNS is untapered and logarithmic; GF, TEF and IEEE are
untapered and linear. The two axes are independent, and our earlier reading ---
that decode cost is set by whether the regime is scanned
(Theorem~\ref{thm:scan}) --- describes only the field-extraction term. It remains
correct within that term, and Table~\ref{tab:realcodecs} shows the term is not
where the money is for a logarithmic format.


\begin{theorem}[The logarithmic decode costs by target precision, not by source width]
\label{thm:targetprec}
A $2^{k}$-entry table for $2^{f}$ with a degree-$d$ Taylor correction delivers
\[
P(k,d) \;=\; (d+1)\Bigl(k + \log_2\tfrac{1}{\ln 2}\Bigr) + \log_2 (d+1)!
\]
bits of relative accuracy, so the table is $P \cdot 2^{k}$ bits with
$k = \frac{P - \log_2 (d+1)!}{d+1} - \log_2\frac{1}{\ln 2}$: exponential in the
target precision, reducible only by raising the polynomial degree, which trades
memory for multipliers. The width of the source format does not enter.
\end{theorem}

\begin{proof}
The residual after the table lookup is $x = f_{\mathrm{lo}}\ln 2$ with
$f_{\mathrm{lo}} < 2^{-k}$, and the degree-$d$ Taylor remainder of $e^{x}$ is
$x^{d+1}/(d+1)!$. Taking $-\log_2$ gives the stated $P$. \qed
\end{proof}

Evaluated against direct simulation of the scheme --- the table, the residual and
the polynomial, in exact arithmetic --- the formula holds to $\pm 0.1$ bit at every
point of a $k \in \{4,8,12,16,20\}$ by $d \in \{1,2,3,4\}$ grid, up to the
$51$-bit resolution of the measurement itself.

Theorem~\ref{thm:targetprec} is falsifiable against the suite and survives. The
takum32 decoder carries $d = 2$ and $k = 16$, and $48/3 = 16$ exactly, where 48 is
the table's output width. If the cost were set by the source width, the takum64
decoder would need a larger table; if it is set by the target, it would need the
same one. It reads \texttt{takum32\_2frac.mem} --- \emph{the same file} --- and
emits binary32, as do the takum8, takum16 and takum64 decoders. One table serves
four source widths because all four target one datapath.

\begin{corollary}[The trade curve, and what each target actually costs]
\label{cor:tradecurve}
Theorem~\ref{thm:targetprec} fixes a curve rather than a point, and the cheapest
point on it that fits a given device is a computable quantity. Pricing the $d$
multiplies at $2.4455 P^{2}$ LUTs --- the quadratic term of our own measured area
law, Theorem~\ref{thm:alpha} --- and requiring the table to fit this device's
$365$ RAMB36 tiles:
\begin{center}
\begin{tabular}{llrr}
\toprule
target & cheapest feasible point & block RAM & LUTs \\
\midrule
binary32 & $d = 2$, $k = 15$ & $43$ tiles & $11{,}269$ \\
binary64 & $d = 5$, $k = 16$ & $187$ tiles & $134{,}808$ \\
\bottomrule
\end{tabular}
\end{center}
An XC7A200T has $134{,}600$ LUTs, so decoding a logarithmic format into a
binary64 datapath costs approximately the entire device.
\end{corollary}

We had previously written that binary64 decode ``is not a design'' at roughly
$3.6$\,Tbit. That figure is the $d = 2$ corner of the curve and taking it for the
whole curve was our error: at $d = 5$ the table is $6.9$\,Mbit and the design
exists. The defensible statement is Corollary~\ref{cor:tradecurve} --- it costs a
device --- and not that it cannot be built.

The measured \texttt{takum32\_decode} sits one $k$ above the binary32 minimum, at
$k = 16$ against $15$, which is a $2\times$ memory margin and ordinary engineering
practice rather than waste. A linear value law needs neither the table nor the
polynomial at any target precision, because the significand is already in the
datapath's representation.

\begin{corollary}[What TEF's exponent path costs]
\label{cor:zeropath}
TEF's value law is linear and its exponent field is fixed, so its exponent path
is wires: no scan, no table, no iteration. Against posit32 that is $480$ LUTs not
spent; against takum32 it is a table of $3.15$\,Mbit not instantiated. This is a
larger and more robust statement than the accuracy result, and unlike it,
Theorem~\ref{thm:floor} places no range bound on it.
\end{corollary}

We flag one asymmetry rather than let it pass. A logarithmic value law buys
something real in exchange: multiplication becomes addition. The comparison above
is of \emph{decoders} --- the cost of getting a value into a binary datapath ---
and a design that stays in the logarithmic domain never pays it. That is the
regime in which LNS accelerators are built, and against those the argument of
Corollary~\ref{cor:zeropath} does not apply at all.


\section{Is there an optimal format, and in what sense}
\label{sec:optimal}

The results above are individually bounding: each says what a format cannot have,
or what a claim may not be extended to. Read together they are stronger than that.
They reduce the design space to four axes, three of which admit unconditional
answers, so that once three inputs are named the format is determined rather than
chosen.

\begin{theorem}[The binary scale is the unique implementable optimum]
\label{thm:radixopt}
Let a format scale by $r^{e}$. Scaling is a shift --- and by
Theorem~\ref{thm:valuelaw} free of any polynomial or table --- if and only if
$r = 2^{k}$ for integer $k \geq 1$. On that set $\kappa(r) = (r-1)^2/(r\ln r)$ is
strictly increasing, taking $0.7213$, $1.6230$, $2.9455$ and $5.0720$ at
$r = 2, 4, 8, 16$. Hence $r = 2$ minimises the expected relative error among all
radices whose scaling is implementable without a transcendental, and it does so
uniquely.
\end{theorem}

\begin{figure}[t]\centering
\includegraphics[width=\linewidth]{tef_radix.pdf}
\caption{$\kappa(r)$ against the scale radix. The unconstrained minimiser sits
near $r = 1.9$ and is unreachable: only $r = 2^{k}$ scales by a shift, and on that
set $\kappa$ is strictly increasing.}
\label{fig:radix}
\end{figure}

The claim is verified by enumeration rather than resting on the monotonicity of
$\kappa$ alone. Sweeping $r$ continuously over $[1.1, 16]$ with $E$ integral and
$M$ taking the remaining width, at eight combinations of width and required range
from $(P,R) = (8, 24)$ to $(32, 5925)$, the best implementable radix is $r = 2$ in
every case, and the penalty against the unreachable optimum is $1.09\times$ at all
eight --- a constant, as it must be, since the two differ only by
$\kappa(1.9)/\kappa(2)$. The unconstrained minimiser is not $2$: it is near
$1.9$, worth $7.9\%$ in error. That gain is unavailable, because any
$r \neq 2^{k}$ makes $r^{e}$ a transcendental evaluation, and
Corollary~\ref{cor:tradecurve} prices those at $84$ RAMB36 tiles --- $23\%$ of
this device --- for a binary32 target. Trading $23\%$ of a chip for $7.9\%$ of a
mantissa bit is not a trade anyone makes. Theorem~\ref{thm:radixopt} is therefore
the interesting kind of optimality result: the optimum is not the extremum, it is
the best point that can be built.

\begin{corollary}[Three inputs determine the format]
\label{cor:determined}
The design space has four axes, and only one of them is a free choice.
\begin{enumerate}
\item \emph{Radix of the scale.} Fixed at $2$ by Theorem~\ref{thm:radixopt},
unconditionally.
\item \emph{Radix of the exponent's encoding.} Fixed by the fabric.
Theorem~\ref{thm:nofree} shows a ternary exponent never spans more values per bit
than a binary one when packed into bits, so it pays on a ternary fabric and
nowhere else.
\item \emph{Fixed fields or a taper.} A genuine choice, and by
Theorem~\ref{thm:dichotomy} exactly two options exist: bounded range at constant
precision, or unbounded range at precision decaying no more gently than one bit
per doubling (Theorem~\ref{thm:floor}).
\item \emph{The split between $E$ and $M$.} Determined once a required range is
named, by Theorem~\ref{thm:law} --- and \emph{only} then, since the error depends
on $M$ alone and the range on $E$ alone.
\end{enumerate}
Name the fabric, the required range, and which side of the dichotomy the workload
sits on, and the format follows.
\end{corollary}

So the question ``is there an ideal format'' has an answer, provided it is asked
with its three arguments supplied. Asked without them it is not underdetermined
but ill-posed: Corollary~\ref{cor:pair} already says a budget cannot be
recommended without a workload, and Theorem~\ref{thm:dichotomy} says no single
format can serve both sides.

\begin{corollary}[What this makes TEF]
\label{cor:whatitmakes}
For the arguments (ternary fabric; a range the workload actually visits; fixed
fields), Corollary~\ref{cor:determined} returns a binary-radix float with a
balanced-ternary exponent field and a uniform mantissa split by
$M = N - 1 - E_t$. That is TEF, and it is forced rather than designed. It follows
that TEF is optimal exactly on those arguments and carries no claim off them ---
in particular none on a binary fabric, where axis 2 returns a binary exponent
instead, and none against a tapered format outside TEF's own range, where
Theorem~\ref{thm:floor} forbids it.
\end{corollary}

\section{Limitations}

\paragraph{The fabric this format is optimal on does not exist to buy.}
Theorem~\ref{thm:nofree} confines the ternary encoding's benefit to a ternary
fabric, and no such process is commercially available. Every hardware figure in
this paper is therefore measured on a binary FPGA, where the ternary exponent
neither gains nor loses --- the enumeration in Section~\ref{sec:whyternary} shows
it ties a binary encoding exactly at the widths tested. The claims that survive on
silicon you can buy are the fixed-field ones: no regime codec, no exponential to
evaluate. The ternary claim is architectural and is labelled as such throughout.
\label{sec:limits}

\begin{enumerate}[leftmargin=*]
\item \textbf{The comparison is against takum16, not tekum16.} The oracle
labelled tekum decodes all $65{,}536$ 16-bit codes identically to the takum
oracle. We report the comparison as against takum and note that a genuine tekum
comparison remains to be done.

\item \textbf{The reference oracle carried a sign defect, now fixed.} Round-tripping
$1.5$ through the parameterised oracle returned $-1.5$ at mantissa widths of 21
and 25 bits: the sign bit sat at a fixed position while the exponent mask took
more bits than the field holds, so the two overlapped. Both are now derived from
the format parameters, and TEF32's round-trip error falls from
$5.4\mathrm{e}{-1}$ --- what a systematically inverted sign averages to --- to
$5.3\mathrm{e}{-9}$.

We record how it survived, because the lesson outlives the bug. The ladder's
check was add/mul commutativity, and an inverted sign survives on both sides of
$a+b=b+a$. The check was real and blind to the class. A round-trip assertion sees
it, costs one line per probe, and now runs at all nine rungs. tekum~\cite{tekum} is a descendant of takum~\cite{takum} adapted for balanced
ternary. Our
oracle reconstructs it from takum's field scheme; the full per-trit
specification requires the tekum paper. The ratios are as good as that model.

\item \textbf{No head-to-head in hardware.} takum's codec RTL is public and
written in VHDL~\cite{takumrtl}. We did not synthesise it alongside TEF, so the cost figures are
TEF's own. What is visible in that source, and which we report as structure
rather than measurement: the 16-bit codec instantiates a 725-line generated
leading-zero-counter and barrel shifter --- the regime decode a fixed-field
format does not have.

\item \textbf{The range is bounded} at $\pm40$ in powers of two, roughly $\pm12$
decades, where a regime field is unbounded. Fixed fields buy the cheap datapath
and the uniform mantissa; range is the price, and workloads that need more should
raise the trit budget or use a tapered format.

\item \textbf{Accuracy above TEF32 is limited by the measuring instrument.}
Once the oracle defect was fixed, TEF32 measures $5.3\mathrm{e}{-9}$, flat
across the workload, which is what 25 mantissa bits should give. TEF64 carries
52 mantissa bits, which is exactly the significand of the IEEE double our metric
is computed in, so the figure it returns describes the metric rather than the
format. Measuring the upper rungs requires exact rationals throughout, as the
oracle itself already uses.

\item \textbf{Six of nine rungs are measured in hardware, and one of those
measurements corrects us.} TEF4 through TEF64 fit the device with hard multipliers
disabled, TEF64 at $7{,}479$ LUTs and $48.20$\,MHz. We previously reported TEF128
as failing to route. It does not fail: with the device's own DSP48 blocks enabled
it closes at $4{,}869$ LUTs, $56$ DSP48 and $103.39$\,MHz, comfortably and at more
than twice TEF64's fabric-only frequency. What fails is TEF128 \emph{under our own}
\texttt{-nodsp} \emph{constraint}, which we impose everywhere so the figures
describe fabric alone and are not confounded by how many hard multipliers a part
happens to carry. That is the right policy for comparing formats; reporting its
consequence as a property of the format was not. The two are different
measurements of different questions. TEF256 upwards exceed a single Artix-7
either way. Table~\ref{tab:spec} is the
specification and the oracle's coverage; Table~\ref{tab:ladder} is what was
placed and routed. We keep them apart on purpose.

\item \textbf{One device family.} Artix-7 on an open flow. Not multi-corner
characterisation; an ASIC mapping will differ.

\item \textbf{The ternary-fabric argument is architectural.} No ternary process
exists to synthesise on. Section~\ref{sec:hardware} is a binary-fabric result;
the ternary claim is reasoning about regime decode and native exponent addition
and should be read as such.
\end{enumerate}

\section*{Reproducibility}

The reference oracle, the RTL, the equivalence testbenches and the synthesis
commands are public. The toolchain is Yosys 0.65, nextpnr-xilinx 1743d0f,
Icarus Verilog 13.0 and Python 3.14 --- all open-source, none requiring a
licence, which is the point: a result nobody can reproduce without buying
something is not a public result.

\section*{Disclosure}

An AI coding assistant was used in preparing the software, the measurements'
tooling and this manuscript. Every figure reported here was produced by running
the named tools on hardware in the author's possession; no measurement in this
paper is estimated, and where a figure is derived rather than observed it is
labelled.

\begin{thebibliography}{9}

\bibitem{takum}
L.~Hunhold.
\newblock Takum arithmetic.
\newblock arXiv:2404.18603.

\bibitem{tekum}
L.~Hunhold.
\newblock Tekum: balanced-ternary tapered-precision arithmetic.
\newblock arXiv:2512.10964.

\bibitem{takumrtl}
L.~Hunhold.
\newblock Design and implementation of a takum arithmetic hardware codec in VHDL.
\newblock arXiv:2408.10594. Source: \url{https://github.com/takum-arithmetic/Takum-Codec-RTL}.

\bibitem{goldenfloat}
D.~Vasilev.
\newblock GoldenFloat: a $\varphi$-derived static-split floating-point family
from GF4 to GF1024 with a Lucas-exact integer identity.
\newblock arXiv:2606.05017v3, 2026.

\bibitem{catalogue}
D.~Vasilev.
\newblock An 83-format numeric catalog with bit-exact conformance vectors: a
vendor-neutral reference for FP8, BF16, MXFP4, and microscaling formats.
\newblock arXiv:2606.09686v2, 2026.

\bibitem{gustafson}
J.~L.~Gustafson and I.~Yonemoto.
\newblock Beating floating point at its own game: posit arithmetic.
\newblock \emph{Supercomputing Frontiers and Innovations}, 4(2), 2017.
\newblock Posit Standard (2022) fixes $es=2$ at all precisions.

\bibitem{yosys}
C.~Wolf.
\newblock Yosys open synthesis suite.
\newblock \url{https://yosyshq.net/yosys/}.

\bibitem{nextpnr}
D.~Shah et al.
\newblock nextpnr: a portable FPGA place-and-route tool.
\newblock \url{https://github.com/YosysHQ/nextpnr}; Xilinx target via Project X-Ray.

\bibitem{hayes2001} B.~Hayes, ``Third Base,'' \emph{American Scientist}, vol.~89, no.~6, pp.~490--494, 2001.
\bibitem{radixeconomy} Radix economy, optimal radix choice. \url{https://en.wikipedia.org/wiki/Radix_economy}
\bibitem{ternary27} Ternary27: A Balanced Ternary Floating Point Format, Version~3.1. \url{https://hackaday.io/project/183153}
\bibitem{goldberg1991} D.~Goldberg, ``What Every Computer Scientist Should Know About Floating-Point Arithmetic,'' \emph{ACM Computing Surveys}, vol.~23, no.~1, pp.~5--48, 1991.
\bibitem{ibmhfp} IBM hexadecimal floating-point. \url{https://en.wikipedia.org/wiki/IBM_hexadecimal_floating-point}
\bibitem{elias1975} P.~Elias, ``Universal codeword sets and representations of the integers,'' \emph{IEEE Trans. Inform. Theory}, vol.~21, no.~2, pp.~194--203, 1975.
\bibitem{positgap} Closing the Gap Between Float and Posit Hardware Efficiency, arXiv:2603.01615.
\bibitem{mlformats} Performance-Efficiency Trade-off of Low-Precision Numerical Formats in Deep Neural Networks, arXiv:1903.10584.
\end{thebibliography}

\end{document}
